\documentclass[11pt]{article}
\usepackage[a4paper,margin=28mm]{geometry}
\usepackage{amsmath,amssymb}
\usepackage{hyperref}
\setlength{\emergencystretch}{2em}
\title{Zonotope faces are not counted by arrangement regions}
\author{Ziang Ni (\texttt{ziangni-sys})}
\date{10 October 2026}
\begin{document}
\maketitle
\noindent\textbf{Disproof of conjecture 00000007132.}
The conjecture equates the face count of a zonotope with the region
count of its hyperplane arrangement. The square already refutes this
for total faces and even for nonempty proper faces.

Let $e_1=(1,0)$ and $e_2=(0,1)$ in $\mathbb R^2$. The zonotope
\[
 Z=[-e_1,e_1]+[-e_2,e_2]=[-1,1]^2
\]
is the Minkowski sum of two actual line segments: each point $(x,y)$
is the sum $(x,0)+(0,y)$, and these sums are precisely the points
satisfying $-1\leq x,y\leq1$.

The associated central hyperplane arrangement consists of
$e_1^\perp=\{x=0\}$ and $e_2^\perp=\{y=0\}$. Its complement has four
regions, namely the four open quadrants. Each quadrant is convex and
nonempty, hence connected. A connected subset of the complement
cannot contain points with opposite signs in either coordinate:
a coordinate projection is continuous, and the intermediate value
theorem would give a point with coordinate zero. Consequently the
four quadrants are exactly the connected components, and the region
count is $4$.

The nonempty proper faces of $Z$ are its four vertices and its four
edges, for a total of $8$. These are genuine exposed faces: a face is
the set of maximizers of a linear functional $ax+by$ on $Z$.
If both coefficients are nonzero, the unique maximizing vertex is
$(\operatorname{sgn}a,\operatorname{sgn}b)$. If exactly one coefficient
is zero, the maximizers form the corresponding edge. If both are
zero, the face is $Z$ itself. This also classifies all exposed faces
of this polytope. The usual total nonempty face count is $9$; counting
the empty face gives $10$. None of $8,9,10$ equals $4$.

For the contradiction it suffices to exhibit five distinct nonempty
proper faces: the right, left, top, and bottom edges, together with
$(1,1)$. Their exposing coefficient pairs are respectively
\[
 (1,0),\quad(-1,0),\quad(0,1),\quad(0,-1),\quad(1,1).
\]
All are different and are strict subsets of the square.

\paragraph{The valid related theorem.}
The familiar region correspondence counts \emph{vertices} of the
zonotope, rather than all its faces. Here both the vertex count and
the arrangement region count are $4$. We do not refute that theorem.
The characteristic polynomial of the two-axis arrangement is
$(q-1)^2$, whose evaluation at $q=-1$ gives $4$ regions; it does not
give the total face count claimed in the source.

\paragraph{Lean coverage.}
The project defines the actual square and its generating coordinate
segments and proves the Minkowski-sum equality. It uses the standard
argmax definition of exposed faces of real linear functionals, proves
five explicit exposing identities, injectivity of the five-face
family, and properness and nonemptiness. It defines the four actual
quadrant sets, proves each connected, and applies the intermediate
value theorem to prohibit axis crossings by any connected subset
of the complement. The accompanying explanation identifies these
sets with the four regions. Lean 4.19.0 and Mathlib are pinned, and
printed audits use only standard logical axioms.
\end{document}
